\begin{pastexam}[label={px:nccu-115-dsa-07}]{政大資科所 115 學年度 第 7 題}
（單選）

(5\%) Which of the following statements regarding the characteristics of several common shortest path
algorithms is incorrect?
\begin{enumerate}[label=(\Alph*)]
\item Dijkstra's algorithm: It is applicable to graphs where all edge weights are non-negative, and it is
typically implemented using a priority queue to achieve a time complexity of $O((|V|+|E|)\lg|V|)$.
\item Bellman-Ford algorithm: It can handle graphs with negative edge weights, but its time complexity is
$O(|V||E|)$, making it slower than Dijkstra's Algorithm.
\item Floyd-Warshall algorithm: It is a single-source shortest path algorithm, suitable for dense graphs,
and solves the problem using the concept of Dynamic Programming.
\item Negative weight cycles: If a graph contains a negative weight cycle, the shortest path is undefined,
and the Bellman-Ford Algorithm can be used to detect the existence of such a cycle.
\end{enumerate}
\end{pastexam}

\begin{pastsoln}
答案是 (C)。

\begin{itemize}
\item (A) 正確。 Dijkstra 的正確性要求邊權非負（資結書第~25.2~節；有負邊時的反例在第~25.2.5~節）。
用二元堆積當優先佇列，$|V|$ 次取出最小值、最多 $|E|$ 次減少鍵值，每次 $O(\lg|V|)$，總共
$O((|V| + |E|)\lg|V|)$。
\item (B) 正確。 Bellman--Ford（資結書第~25.3~節）不怕負邊，做 $|V| - 1$ 輪、每輪鬆弛全部 $|E|$ 條邊，
$O(|V||E|)$。 $|E| \ge |V| - 1$ 的連通圖上，$|V||E|$ 比 $(|V| + |E|)\lg|V|$ 大，所以比 Dijkstra 慢。
\item (C) 錯誤，就是要選的。 Floyd--Warshall（資結書第~25.6~節）是\textbf{全點對}最短路徑演算法，
一次算出每一對點的距離，$\Theta(|V|^3)$，和邊數無關，所以適合稠密圖；它的確是動態規劃
（子問題是「中間點只能用前 $k$ 個點」）。 錯的是「single-source」這幾個字。
\item (D) 正確。 從起點走得到一個負環時，繞環一圈距離就變小，可以無限變小，最短路徑沒有定義。
Bellman--Ford 做完 $|V| - 1$ 輪後再檢查一輪，若還有邊能鬆弛，就表示有從起點走得到的負環
（資結書第~25.3.4~節）。
\end{itemize}
\end{pastsoln}
